\documentclass[11pt]{article}
\usepackage[T1]{fontenc}
\usepackage{lmodern,amsmath,amssymb,amsthm}
\usepackage[margin=1in]{geometry}
\usepackage[colorlinks=true,linkcolor=blue,citecolor=blue,urlcolor=blue]{hyperref}
\newtheorem{theorem}{Theorem}
\newtheorem{proposition}[theorem]{Proposition}
\newtheorem{lemma}[theorem]{Lemma}
\newcommand{\R}{\mathbb R}
\newcommand{\Sph}{\mathbb S^2}
\newcommand{\dd}{\,d}
\DeclareMathOperator{\sech}{sech}
\title{The Mercedes-star partition of a disk}
\author{Kaiwen Jin\\\small\href{mailto:yui.ui.siki@gmail.com}{yui.ui.siki@gmail.com}}
\date{October 12, 2026}
\begin{document}
\maketitle
\begin{abstract}
Among triples of pairwise disjoint nonempty connected open subsets of the
unit disk, the three sectors of opening $2\pi/3$ minimize the largest first
Dirichlet eigenvalue. The minimum is $j_{3/2,1}^2$, where $j_{3/2,1}$ is
the first positive zero of the Bessel function $J_{3/2}$.
The proof transfers any triple with smaller energy to a triple on the unit
sphere with energy below its known sharp minimum $15/4$. A common radial
change of variables matches the angular terms after division by positive
reference profiles. A Wronskian comparison shows that the remaining radial
term has the required sign. The argument imposes no condition on the cells'
boundary regularity, topology, symmetry, or coverage of the disk.
\end{abstract}

\section{The minimax problem}
Spectral partition problems balance the first Dirichlet eigenvalues of
several disjoint cells. For three cells in a disk, the natural candidate is
the Mercedes star: three equal sectors meeting at the center.
Let $B=\{x\in\R^2:|x|<1\}$. For a nonempty open subset $D$ of $B$, write
\begin{equation}\label{eq:rayleigh}
 \lambda_1(D)=\inf_{0\ne u\in C_c^\infty(D;\R)}
                  \frac{\int_D|\nabla u|^2}{\int_Du^2}.
\end{equation}
This also equals the infimum over $H_0^1(D)$, defined as the closure of
$C_c^\infty(D)$ in $H^1(D)$: convergence in that norm passes the energy
inequality to the closure, and the reverse comparison follows from inclusion.
This convention is valid for arbitrary open sets.

Berkolaiko, Canzani, Cox and Marzuola prove minimality of the Mercedes star
among partitions whose boundary contains a curve from the origin to the
boundary of the disk \cite[Theorem~1.8]{BCCM}. The unrestricted class also
allows interfaces outside that comparison. On the unit sphere, in contrast,
the three equal lunes are known to minimize the same objective
\cite[Theorem~1.1 and (1.9)]{HHT}. The arbitrary-open-cell formulation is
stated by Soave and Terracini \cite[Definition~1 and Theorem~1.10]{ST}.
We use that spherical result to obtain the unrestricted disk comparison.

\begin{theorem}\label{thm:main}
For every triple $(D_1,D_2,D_3)$ of pairwise disjoint nonempty connected open
subsets of $B$,
\[
 \max_{1\le i\le3}\lambda_1(D_i)\ge j_{3/2,1}^2.
\]
The three sectors separated by radii at angles $2\pi/3$ attain equality.
Thus the infimum over all such triples is $j_{3/2,1}^2$.
Here $j_{3/2,1}$ is the first positive zero of the Bessel function $J_{3/2}$.
\end{theorem}

The comparison acts on arbitrary compactly supported test functions.
We divide out the positive radial profile of an equal sector, and choose a
radial map so that its angular energy agrees with the corresponding
spherical energy. The radial energy is then larger. Consequently, three
disk cells with eigenvalues below the proposed value would yield three
spherical cells below their sharp value. A logarithmic cutoff handles tests
that meet the center. We prove the radial comparison in Section~2, the
energy identity in Section~3, and the theorem in Section~4.

We record the external spectral result in the precise form used below.
The sphere has radius one and its induced metric; its eigenvalues are
defined by \eqref{eq:rayleigh} with spherical gradient and area measure.
\begin{proposition}[Spherical three-partition theorem]\label{prop:sphere}
For every triple of pairwise disjoint nonempty connected open subsets
$\Omega_i$ of $\Sph$, one has
$\max_i\lambda_1(\Omega_i)\ge15/4$.
\end{proposition}
This is the $N=3$ case of \cite[Theorem~1.10]{ST}, with the partition class
of Definition~1 and the objective in equation~(10) of that paper's
arXiv v2. It requires neither a covering nor regular boundaries.

\section{A radial comparison}
Put $\nu=3/2$ and $E=j_{3/2,1}^2$. It is useful to write the positive
reference profile without Bessel notation. Let $k$ be the first positive
zero of $f(x)=\sin x-x\cos x$. Since $f'=x\sin x$, it is positive on
$(0,\pi]$ and strictly decreases from $\pi$ to $-1$ on
$(\pi,3\pi/2)$. Hence $k\in(\pi,3\pi/2)$.
The identities for the spherical Bessel function
\[
 j_1(z)=\sqrt{\frac{\pi}{2z}}J_{3/2}(z)
       =\frac{\sin z}{z^2}-\frac{\cos z}{z}
 \quad(z>0)
\]
\cite[10.47.3 and 10.49.3]{DLMF} identify $k=j_{3/2,1}$.
Define
\begin{equation}\label{eq:q}
 q(r)=r^{-3/2}\{\sin(kr)-kr\cos(kr)\},\qquad 0<r<1.
\end{equation}
It is positive and solves
\begin{equation}\label{eq:diskode}
 -q''-\frac{q'}r+\frac{\nu^2}{r^2}q=Eq.
\end{equation}
Indeed, for $g(r)=f(kr)$, $g'=k^2r\sin(kr)$ and
$g''=k^2\sin(kr)+k^3r\cos(kr)$; the left side is
$r^{-3/2}(2g'/r-g'')=k^2q$. At the endpoints,
\begin{equation}\label{eq:endpoints}
 \begin{aligned}
 q(r)&=a r^{3/2}+O(r^{7/2}),&
 q'(r)&=\tfrac32a r^{1/2}+O(r^{5/2}),\\
 a&=k^3/3>0,& q(1)&=0,\qquad q'(1)=k^2\sin k<0.
 \end{aligned}
\end{equation}

The spherical profile will be $\sin^{3/2}\theta$. To match the angular
energy after dividing by these profiles, set
\begin{equation}\label{eq:massmap}
 \begin{gathered}
 T=\int_0^1q(r)^2\frac{\dd r}{r},\qquad C=\frac{2T}{\pi},\\
 \int_0^{\theta(r)}\sin^2 a\dd a
       =\frac1C\int_0^r q(t)^2\frac{\dd t}{t}.
 \end{gathered}
\end{equation}
Here $0<T<\infty$ by \eqref{eq:endpoints}.
The cumulative functions increase strictly from $0$ to $\pi/2$;
therefore $\theta:(0,1)\to(0,\pi)$ is a smooth increasing bijection,
with endpoint limits $0$ and $\pi$. In particular,
\begin{equation}\label{eq:thetader}
 \theta'=\frac{q^2}{Cr\sin^2\theta}.
\end{equation}

\begin{lemma}\label{lem:radial}
For every $0<r<1$,
\begin{equation}\label{eq:h}
 q(r)^2>C\sin^3\theta(r),\qquad
 h(r):=\frac{r\theta'(r)}{\sin\theta(r)}>1.
\end{equation}
\end{lemma}
\begin{proof}
Use logarithmic radius: $Q(t)=q(e^t)$ on $t<0$, extended by zero on
$t\ge0$. For each real $s$, compare it with
$Z_s(t)=\sqrt C\,\sech^{3/2}(t-s)$ on $\R$.
The substitution $a=2\arctan e^{t-s}$ gives
$a'=\sin a=\sech(t-s)$, and hence
\[
 \int_\R Q^2\dd t=T=\int_\R Z_s^2\dd t.
\]
On $t<0$, \eqref{eq:diskode} and differentiation give
\[
 Q''=(\nu^2-Ee^{2t})Q,\qquad
 Z_s''=(\nu^2-\nu(\nu+1)\sech^2(t-s))Z_s.
\]
The Wronskian $W=Q'Z_s-QZ_s'$ therefore satisfies
\begin{equation}\label{eq:wronskian}
 W'=e^{2t}b_s(t)QZ_s,\qquad
 b_s(t)=\frac{15e^{-2s}}{(1+e^{2(t-s)})^2}-E.
\end{equation}
The coefficient $b_s$ decreases strictly, and
$b_s(0)\le15/4-E<0$.
Also $W(-\infty)=0$, because both profiles and their derivatives are
$O(e^{3t/2})$, while $W(0^-)=q'(1)Z_s(0)<0$.
If $b_s$ is negative throughout, then $W<0$. Otherwise $W$ first increases
from zero and then strictly decreases, crossing zero once before $t=0$.
Since $(Q/Z_s)'=W/Z_s^2$, the ratio is strictly decreasing, or increases
and then decreases. Its zero extension adds no peak.

It follows that $I=\{Q^2>Z_s^2\}$ is a single open interval. It is
nonempty, since the two total masses agree and $Q^2-Z_s^2<0$ on
$[0,\infty)$. Its right endpoint $b$ is less than zero.
Write $H(t)=\int_{-\infty}^t(Q^2-Z_s^2)\dd a$.
If $I=(-\infty,b)$, then $H(t)>0$ at every finite $t$: before $b$ it is
an integral of a positive function, and after $b$ it is minus the remaining
negative tail integral. If $I=(a,b)$ with finite $a$, then $H<0$ up to $a$,
$H$ is strictly increasing on $(a,b)$, and
\[
 H(b)=-\int_b^\infty(Q^2-Z_s^2)\dd t>0.
\]
Outside $I$ the density difference is strictly negative except at its
endpoints; strict monotonicity of the ratio excludes an equality interval.
Thus any finite zero of $H$ lies strictly inside $I$. We have proved
\begin{equation}\label{eq:crossing}
 \int_{-\infty}^{t_0}Q^2\dd t=
 \int_{-\infty}^{t_0}Z_s^2\dd t
 \quad\Longrightarrow\quad Q(t_0)^2>Z_s(t_0)^2
 \qquad(t_0<0).
\end{equation}

For a fixed $r$, choose $t_0=\log r$ and
$s=t_0-\log\tan(\theta(r)/2)$.
The substitution defining $a$ shows that \eqref{eq:massmap} is exactly
the cumulative equality in \eqref{eq:crossing}.
Since $\sech(t_0-s)=\sin\theta(r)$, this gives the first inequality in
\eqref{eq:h}; the second follows from \eqref{eq:thetader}.
\end{proof}

\section{Comparison of quadratic forms}
The radial map
\begin{equation}\label{eq:map}
 \Phi(r,\phi)=(\sin\theta(r)\cos\phi,
               \sin\theta(r)\sin\phi,\cos\theta(r)),
 \qquad\Phi(0)=(0,0,1),
\end{equation}
is a homeomorphism from $B$ onto the sphere minus its south pole.
Indeed, the radial functions are inverse continuous increasing bijections,
and the angular circles collapse continuously at the origin and north pole.
Away from these two points, it is a smooth diffeomorphism.
Its image is open in the full sphere, so it preserves openness,
connectedness, nonemptiness and disjointness of cells.

Let $u$ be a real compact smooth test on $B$ whose support avoids the origin.
Set
\[
 u=qv,\qquad V(\theta,\phi)=v(r(\theta),\phi),\qquad
 w(\theta,\phi)=\sin^\nu\theta\,V(\theta,\phi).
\]
Extend $w$ by zero near the two poles. It is smooth there because its
support stays away from both poles, and it is nonzero whenever $u$ is.
Define the shifted forms
\[
 G_B(u)=\int_B|\nabla u|^2-E\int_Bu^2,\qquad
 G_S(w)=\int_{\Sph}|\nabla_Sw|^2-\tfrac{15}{4}\int_{\Sph}w^2.
\]
Expanding $u_r=q'v+qv_r$ and integrating the mixed term gives
\begin{equation}\label{eq:gb}
 G_B(u)=\int_0^{2\pi}\int_0^1
 \left[rq^2v_r^2+\frac{q^2}{r}(v_\phi^2-\nu^2v^2)\right]\dd r\dd\phi.
\end{equation}
For clarity, the coefficient of $v^2$ after that integration is
$r q'^2-Erq^2-(rqq')'=-q(rq')'-Erq^2=-\nu^2q^2/r$.
All boundary terms vanish by compact support.

On the sphere the area measure is $\sin\theta\dd\theta\dd\phi$ and
$|\nabla_Sw|^2=w_\theta^2+w_\phi^2/\sin^2\theta$.
The profile $p(\theta)=\sin^\nu\theta$ satisfies
\[
 -p''-\cot\theta\,p'+\frac{\nu^2}{\sin^2\theta}p
       =\nu(\nu+1)p=\tfrac{15}{4}p.
\]
This follows from $p'/p=\nu\cot\theta$ and
$\csc^2\theta-\cot^2\theta=1$.
The corresponding mixed-term integration yields
\begin{equation}\label{eq:gs}
 G_S(w)=\int_0^{2\pi}\int_0^\pi
 \left[\sin^4\theta\,V_\theta^2+
       \sin^2\theta\,(V_\phi^2-\nu^2V^2)\right]\dd\theta\dd\phi.
\end{equation}
There are again no boundary terms, since the support avoids the poles.
By \eqref{eq:thetader} and \eqref{eq:h},
\[
 \frac{q^2}{r\theta'}=C\sin^2\theta,
 \qquad rq^2\theta'=C\sin^4\theta\,h^2.
\]
Changing variables in \eqref{eq:gb} and subtracting $C$ times
\eqref{eq:gs} proves the exact identity
\begin{equation}\label{eq:comparison}
 G_B(u)-CG_S(w)=C\int_0^{2\pi}\int_0^\pi
  \sin^4\theta\,[h(r(\theta))^2-1]V_\theta^2\dd\theta\dd\phi\ge0.
\end{equation}
Thus a negative disk form gives a negative spherical form. In particular,
no comparison of the two unweighted masses is needed.

\section{Proof of Theorem~\ref{thm:main}}
Suppose that an admissible triple had $\lambda_1(D_i)<E$ for all three
indices. Choose nonzero $u_i\in C_c^\infty(D_i;\R)$ with $G_B(u_i)<0$.
We first make their supports avoid the origin without losing these strict
inequalities.
Choose a smooth $\eta:\R\to[0,1]$, zero on $(-\infty,0]$ and one on
$[1,\infty)$, and put
\[
 \chi_\varepsilon(r)=\eta\left(
 \frac{\log r-2\log\varepsilon}{\log(1/\varepsilon)}\right),
 \qquad 0<\varepsilon<1/2.
\]
It vanishes for $r\le\varepsilon^2$ and equals one for $r\ge\varepsilon$.
Its constant extension near zero is smooth. With $M=\|\eta'\|_\infty$,
\[
 \int_B|\nabla\chi_\varepsilon|^2
       \le\frac{2\pi M^2}{\log(1/\varepsilon)}\longrightarrow0.
\]
For a fixed compact smooth $u$, the $L^2$ errors
$(\chi_\varepsilon-1)u$ and
$(\chi_\varepsilon-1)\nabla u$ are bounded by their supremum norms times
$\sqrt\pi\varepsilon$, while
$\|u\nabla\chi_\varepsilon\|_2\le\|u\|_\infty
\|\nabla\chi_\varepsilon\|_2\to0$.
Hence $\chi_\varepsilon u\to u$ in $H^1(B)$.
A common sufficiently small $\varepsilon$ leaves all three products
nonzero with negative $G_B$, still compactly supported in their original
cells. Relabel them $u_i$.

The cells $\Omega_i=\Phi(D_i)$ are pairwise disjoint nonempty connected
open spherical sets. The tests obtained in Section~3 are nonzero members of
$C_c^\infty(\Omega_i)$, with $G_S(w_i)<0$ by \eqref{eq:comparison}.
Thus all three $\lambda_1(\Omega_i)<15/4$, contradicting
Proposition~\ref{prop:sphere}. This proves the universal lower bound.

For the upper bound, let $\ell=2\pi/3$ and consider the sector
$0<r<1$, $0<\phi<\ell$. Its trial function
$U=q(r)\sin(\nu\phi)$ has finite mass and gradient energy by
\eqref{eq:endpoints}. Multiplication of \eqref{eq:diskode} by $rq$ gives
\[
 \int_0^1\left(rq'^2+\nu^2q^2/r\right)\dd r
        =E\int_0^1rq^2\dd r,
\]
since $rqq'\to0$ at both endpoints. Both squared angular factors integrate
to $\ell/2$, so the quotient of $U$ is $E$.

To verify $U\in H_0^1$ explicitly, choose a smooth $\rho$ bounded between
zero and one, zero on $(-\infty,1]$ and one on $[2,\infty)$.
Multiply $U$ by
\[
 \rho(r/\delta)\rho((1-r)/\delta)
 \rho(\phi/\delta)\rho((\ell-\phi)/\delta).
\]
For small $\delta>0$ this is a compact smooth sector test.
The original mass and gradient integrals over the removed strips tend to
zero by integrability. The extra derivative energies are $O(\delta^3)$
at the vertex, since $q=O(r^{3/2})$, and $O(\delta)$ at the circular edge,
since $q=O(1-r)$. At either angular edge the sine vanishes linearly,
and the extra energy is bounded by
\[
 K\left(\int_0^1q^2\frac{\dd r}{r}\right)
    \delta^{-2}\int_0^{2\delta}a^2\dd a=O(\delta).
\]
Here $K$ is independent of $\delta$; the other cutoffs are bounded by one.
The square of the sum of four extra gradient terms is at most four times
the sum of their squares. The tests therefore converge to $U$ in $H^1$.
Their quotients tend to $E$, giving the sector eigenvalue at most $E$.
The three rotated sectors are an admissible triple. Their maximum is at
most $E$ and, by the lower bound, at least $E$. This proves attainment and
the stated exact infimum.\hfill$\square$

\begin{thebibliography}{9}
\bibitem{BCCM}
G. Berkolaiko, Y. Canzani, G. Cox and J. L. Marzuola,
\emph{Homology of spectral minimal partitions},
arXiv:2406.04225v1 (2024).
\url{https://arxiv.org/abs/2406.04225v1}.
\bibitem{HHT}
B. Helffer, T. Hoffmann-Ostenhof and S. Terracini,
\emph{On spectral minimal partitions: the case of the sphere},
arXiv:0903.3326v1 (2009).
\url{https://arxiv.org/abs/0903.3326v1}.
\bibitem{ST}
N. Soave and S. Terracini,
\emph{Liouville theorems and $1$-dimensional symmetry for solutions of an
elliptic system modelling phase separation},
Advances in Mathematics \textbf{279} (2015), 29--66.
\href{https://doi.org/10.1016/j.aim.2015.03.015}{doi:10.1016/j.aim.2015.03.015}.
Source version: \url{https://arxiv.org/abs/1404.7288v2}.
\bibitem{DLMF}
NIST, \emph{Digital Library of Mathematical Functions},
equations 10.47.3 and 10.49.3.
\url{https://dlmf.nist.gov/10.47.E3};
\url{https://dlmf.nist.gov/10.49.E3}.
\end{thebibliography}

\paragraph{AI-use statement.}
OpenAI Codex was used for proof exploration, source checks, symbolic
calculations and drafting this manuscript.
\end{document}
